\documentclass[11pt,a4paper]{article}
\usepackage[utf8]{inputenc}
\usepackage[T1]{fontenc}
\usepackage{geometry}
\geometry{top=1in, bottom=1in, left=1in, right=1in}
\usepackage{hyperref}
\usepackage{booktabs}
\usepackage{amsmath}
\usepackage{listings}
\usepackage{xcolor}
\usepackage{graphicx}
\usepackage{array}
\usepackage{fancyhdr}
\usepackage{titlesec}
\usepackage{longtable}

\hypersetup{
    colorlinks=true,
    linkcolor=purple,
    filecolor=magenta,      
    urlcolor=cyan,
    pdftitle=ADR-001: The Repository Folder name must be portal/,
    pdfpagemode=FullScreen,
}

\definecolor{codegray}{rgb}{0.96,0.96,0.96}
\definecolor{codeborder}{rgb}{0.85,0.85,0.85}
\definecolor{commentgreen}{rgb}{0.12,0.55,0.22}
\definecolor{keywordblue}{rgb}{0.1,0.2,0.8}

\lstset{
    backgroundcolor=\color{codegray},
    basicstyle=\ttfamily\small,
    breakatwhitespace=false,
    breaklines=true,
    captionpos=b,
    commentstyle=\color{commentgreen},
    frame=single,
    rulecolor=\color{codeborder},
    keepspaces=true,
    keywordstyle=\color{keywordblue},
    language=Python,
    showspaces=false,
    showstringspaces=false,
    showtabs=false,
    tabsize=4
}

\pagestyle{fancy}
\fancyhf{}
\rhead{\small ADR-001: The Repository Folder name must be portal/}
\lhead{\small Project Genesis}
\cfoot{\thepage}

\title{\textbf{ADR-001: The Repository Folder name must be portal/}}
\author{Project Genesis Research Repository}
\date{\small Generated: July 2026}

\begin{document}

\maketitle
\thispagestyle{empty}

\newpage
\section*{ADR-001: The Repository Folder name must be \textbackslash{}texttt{portal/}}

\subsection*{Status}
Approved

\subsection*{Context}
The directory structure needs a clear identifier for the web-facing component.

\subsection*{Decision}
We use \textbackslash{}texttt{portal/} instead of \textbackslash{}texttt{website/}, \textbackslash{}texttt{frontend/}, or \textbackslash{}texttt{site/}.

\subsection*{Consequences}
Future desktop client binaries or visualization clients can sit alongside it in the repository under their own subdirectories (e.g., \textbackslash{}texttt{client/}) without architectural confusion.

\end{document}
